\section{Appendix}
\label{sec:appendix}

\subsection{Reproducibility}

The experiments are implemented in the ListT5 project notebooks.
The v2 notebook contains the strategy-aware caching update used for the score-balanced aligned-cache runs.

The key configuration for the preliminary run is:
\begin{itemize}
  \item model: \texttt{Soyoung97/ListT5-base}
  \item datasets: \texttt{nfcorpus}, \texttt{scifact}, \texttt{arguana}, \texttt{scidocs}, \texttt{fiqa}
  \item maximum queries per dataset: 50
  \item maximum input sequence length: 512
  \item listwise group size: 5
  \item output size per listwise call: 2
  \item reranking depth: 10
  \item BM25 top-k comparison: sequential top-40, top-60, top-80, and top-100
  \item grouping comparison: sequential top-100 vs score-balanced top-100
  \item v2 caching check: strategy-aware first-round cache for score-balanced grouping
\end{itemize}

\subsection{Grouping Example}

With 20 candidates and group size 5, sequential grouping forms:
\begin{quote}
\texttt{[1,2,3,4,5]}, \texttt{[6,7,8,9,10]}, \texttt{[11,12,13,14,15]}, \texttt{[16,17,18,19,20]}.
\end{quote}

Score-balanced grouping forms:
\begin{quote}
\texttt{[1,5,9,13,17]}, \texttt{[2,6,10,14,18]}, \texttt{[3,7,11,15,19]}, \texttt{[4,8,12,16,20]}.
\end{quote}

The implementation uses zero-based candidate indices internally, but the paper describes ranks using one-based indexing for readability.
